\section{Full text of routine threats to validity}
\label{sec:appendix_threats}

This appendix carries the full \emph{threat / mitigation / residual risk}
text for the three routine methodological threats summarised in
Table~\ref{tab:threats_routine} in the main body
(Section~\ref{sec:limitations}). The three load-bearing threats
(judge SLM validity, BERTScore construct validity, benchmark coverage
and contamination) remain as prose in the main body.

\subsection{Hardware split during the sweep}
\label{sec:appendix_hardware_split}

\textbf{Threat.} The 22 GPU-days of experimental sweep were split
across Google Colab A100 and T4 hardware. Different cells were
executed on different hardware because A100 availability was
intermittent; the mono stratum (M1--M6) and the initial pilot ran on
A100, the H8/H10/H11/H6/H9 hetero cells ran on T4, and the remaining
cells split across both. Different hardware could in principle produce
different LLM outputs due to different floating-point kernels or
different attention implementations.

\textbf{Mitigation.} The Ollama runtime enforces greedy decoding at
temperature $0.2$ across all models used in this study. All model
weights are quantised and deterministic under the same seed and
prompt. The LLM disk cache is keyed on (\texttt{model\_tag},
\texttt{role\_hint}, \texttt{prompt}, \texttt{temperature},
\texttt{top\_p}, \texttt{top\_k}, \texttt{num\_ctx}, \texttt{max\_tokens})
--- identical parameter tuples produce identical cache hits regardless
of hardware. Cross-hardware verification via re-execution of 30 pilot
rows on both hardware types produced byte-identical LLM outputs.

\textbf{Residual risk.} Rare non-determinism in the Ollama runtime
under GPU memory pressure could produce different output text on a
small number of rows. We estimate this affects fewer than $0.1\%$ of
rows based on the observed hit rate of the cache during resumption
from mid-sweep checkpoints.

\paragraph{Related: mid-sweep NUM\_CTX configuration change.} During
the sweep, a configuration change reduced $\texttt{NUM\_CTX}$ from
$16384$ to $8192$ tokens for the affected models to work around T4
memory constraints. Rows generated before and after the change have
subtly different cache keys, so the G2 per-mutation-operator
regeneration script achieved only $67\%$ cache-recovery coverage on
Path~1 rows. The per-operator table (Table~\ref{tab:per_operator_kill_rate})
is computed only on the $1{,}172$ Path~1 rows for which cache
reconstruction succeeded; the remaining $567$ rows are excluded, not
silently treated as zeros. All other analyses use the aggregate
metric value stored directly in the TSV, which is unaffected by the
cache-key mismatch.

\subsection{Cache reuse across cells and per-path duplicate cells}
\label{sec:appendix_cache_reuse}

\textbf{Threat.} The five-layer checkpointing system (LLM cache,
pytest cache, BERTScore cache, TSV row cache, JSONL decontamination
cache) means that identical calls across cells are served from cache.
This introduces a specific analytical concern: each closure path
exercises only two of the three pipeline stages (Path~1 uses
$L_{\textsf{spec}}$ and $L_{\textsf{test}}$; Path~2 uses
$L_{\textsf{test}}$ and $L_{\textsf{code}}$; Path~3 uses
$L_{\textsf{test}}$ and $L_{\textsf{spec}}$), so cells sharing the
operative stage-pair on a given path can produce identical rows via
cache reuse. We measured the empirical per-group identity rate
directly from the released sweep TSV, counting samples with identical
$(m, J)$ tuples across every cell in a candidate group:

\begin{itemize}
  \item \textbf{Genuine duplicates ($\geq 98\%$ identical rows on shared samples):}
    $\{M6, H4\}$ on Path~1 ($65/65$, $100\%$) and Path~3 ($74/75$, $98.7\%$);
    $\{M6, H1, H5\}$ on Path~2 ($74/75$, $98.7\%$).
    These groups share the operative stage-pair on the affected path
    and were swept entirely within the same cache era; the effective
    per-cell degrees of freedom on the affected path collapse across
    the group.
  \item \textbf{Partial duplicates ($\sim$$50\%$ identical rows):}
    $\{H7, H9\}$ on Path~1 ($33/59$, $56\%$);
    $\{H2, H10\}$ on Path~2 ($37/75$, $49\%$);
    $\{H6, H7\}$ on Path~2 ($38/75$, $51\%$).
    The partial rate is a consequence of the mid-sweep
    $\texttt{NUM\_CTX}$ configuration change disclosed above:
    pre-change rows are cache-identical across the group while
    post-change rows have different cache keys.
  \item \textbf{Notably \emph{not} duplicates:}
    $\{H7, H9\}$ on Path~3 ($0/74$ identical) --- both cells share
    $(L_{\textsf{test}}, L_{\textsf{spec}}) = (\texttt{qwen3.6}, \texttt{qwen-coder})$
    on Path~3 by our indexing convention, but no rows survive with
    byte-identical BERTScore $+$ judge outputs on this path.
    $\{M1, N1\}$ on every path ($9/51 = 17.6\%$ Path~1, $15/75 = 20.0\%$
    Path~2, $3/75 = 4.0\%$ Path~3): although M1 and N1 share
    $L_{\textsf{spec}} = L_{\textsf{test}} = L_{\textsf{code}} = \texttt{llama3.2}$,
    N1 is a word-shuffled-input null control
    (Section~\ref{sec:experiment_setup_doe}) so its per-sample inputs
    differ from M1's by construction --- the shared-stage-pair
    reasoning does not apply. The $\bar{Y}$ scale difference (M1
    $\bar{Y}_1 = 0.577$ vs.\ N1 $\bar{Y}_1 = 0.429$; false-closure
    rate $14.0\%$ vs.\ $26.3\%$) confirms the two cells are
    behaviorally distinct, which is precisely what allows N1 to serve
    as the RQ4 noise floor in Section~\ref{sec:discussion_rq4}.
\end{itemize}

The apparent H7 $=$ H9 byte-identity on Path~1 in the cache-reconstructed
per-operator table (Table~\ref{tab:per_operator_kill_rate}) reflects the
cache-recoverable subset of Path~1 rows (which necessarily excludes the
post-$\texttt{NUM\_CTX}$-change rows where H7 and H9 diverge), not the
released sweep TSV. Readers checking H7 vs.\ H9 across the TSV and the
per-operator table will find the two artifacts differ, and this is why.

\textbf{Mitigation.} We disclose the measured duplication rates
explicitly. In interpreting the per-path ANOVAs
(Section~\ref{sec:results_stage_attribution}) and the mixed-effects
logit (Section~\ref{sec:results_path_cell_interaction}), the effective
per-path degrees of freedom on genuine-duplicate groups collapse across
the group (M6/H4 counts once on Paths~1 and~3; M6/H1/H5 counts once on
Path~2), and partial-duplicate groups count with a scaled-down
effective sample. The qualitative RQ1--RQ3 conclusions (heterogeneity
matches mono on aggregate $\bar{Y}$; test-stage is the significant
per-stage attribution; the phi-4-in-test-stage pathology) do not depend
on between-duplicate contrasts and survive the correction. The
mixed-effects logit is the least affected because it estimates
per-model-per-stage marginal effects rather than per-cell effects;
genuine duplicates on a given path collapse to the same fitted
probability.

\textbf{Residual risk.} Tukey contrasts within the two genuine-duplicate
groups (M6/H4 on Paths~1 and~3; M6/H1/H5 on Path~2) overstate the
effective per-comparison sample size. No such within-group contrast
appears among the 25 significant pairs listed in
Table~\ref{tab:tukey_significant}; the contrasts highlighted in the
text of Section~\ref{sec:results_mono_vs_hetero} (H1--M1, H4--N1,
M1--N2) are all between-group. Between-group contrasts and
within-partial-duplicate contrasts are affected proportionally to the
measured identity rate.

\subsection{Multiple testing}
\label{sec:appendix_multiple_testing}

\textbf{Threat.} The paper reports many statistical tests: the main
ANOVA (Eq.~\ref{eq:anova}), the interaction ANOVA
(Eq.~\ref{eq:anova_interact}), 190 pairwise Tukey comparisons, three
per-path Pearson correlations, per-benchmark ANOVAs, per-operator
effect sizes, capability-preservation rates. Without correction, some
significance would be expected by chance.

\textbf{Mitigation.} The Tukey HSD pairwise test controls family-wise
error at $\alpha = 0.05$ over the 190-pair family via its own
studentized-range distribution (Section~\ref{sec:experiment_setup_stats});
no additional Bonferroni layering is applied. Per-path correlations and
their p-values are reported directly without correction; the significant
ones (Path~2, $r_2 = 0.52$, $p = 1.19 \times 10^{-129}$; per-benchmark
Path~3, both $p \leq 0.001$) are far below the significance threshold
that no correction affects their interpretation, and the aggregate
Path~3 $r_3 = 0.022$, $p = 0.359$ is reported alongside its
Simpson's-paradox interpretation rather than as a null claim. The
interaction ANOVA is a single planned comparison. The main-effect
ANOVAs are also planned comparisons.

\textbf{Residual risk.} Post-hoc per-operator analyses in
Section~\ref{sec:results_per_operator} are not corrected for multiple
testing beyond what the reported effect sizes convey. Effect sizes
$\Delta > 0.05$ are treated as meaningful; smaller effect sizes are
not emphasized.
